\clearpage\begingroup
\HistoricalHead{165}{R}{GAUSS AN SCHUMACHER UND OLBERS.}
\begin{center}
[23.]\par\medskip
\textbf{\textsc{Gauss} an \textsc{Olbers}.\quad Göttingen, 25. Januar 1832.}\par\medskip
\rule{.78\linewidth}{.4pt}\par\smallskip
{\footnotesize [\textsc{Wilhelm Olbers}, Sein Leben und seine Werke II, 2, 1900, S. 581.]}\par\smallskip
\rule{.78\linewidth}{.4pt}
\end{center}
\noindent . . . . . . . . . . . . . . . . . . . . .\par
\noindent\hspace*{\parindent}Sie haben vielleicht vor einiger Zeit in den G. G. A.[*)]\begingroup\renewcommand{\thefootnote}{}\footnotetext{[*) Göttingische Gelehrte Anzeigen, 1831, 196. Stück, S. 1947, siehe den unten folgenden \textit{Anhang}.]}\endgroup\ gesehen, dass der Preis auf verbesserte Photometer nicht hat zuerkannt werden können, wenn gleich zwei Abhandlungen eingegangen waren. Die zweite, deren Verfasser sich übrigens fast mit Gewissheit errathen liess, enthält manches Eigenthümliche, auch einen Vorschlag, der auf einer an sich sehr sinnreichen Idee beruht, die ich aber bei näherer Prüfung nicht stichhaltig gefunden habe. Übrigens kann ich nicht leugnen, dass mir scheint, alle Methoden, die sich darauf gründen, dass man das Licht so lange schwächt, bis es einem gegebenen Auge verschwindet, könnten nur sehr rohe Resultate geben; jedenfalls sollte
\par\clearpage\SourceLeaf{171}\HistoricalHead{166}{L}{PHYSIK.\quad BRIEFWECHSEL.}
\noindent man damit anfangen zu prüfen, welche Übereinstimmung unter sich sie geben, und zwar nicht an 3 oder 6 Messungen, sondern an einigen Hunderten. Ich bin geneigt zu glauben, dass allemahl, wo das kleinste sichtbare Licht die Einheit abgeben soll, Ungewissheiten von 30 oder mehreren p[ro] c[ent] nicht zu vermeiden sein werden.

Ich habe mir eine Methode ausgesonnen, die auf ein anderes Princip sich stützt, wo nemlich das Auge über nichts als über Gleichheit oder Ungleichheit zweier Gegenstände, die es immer nebeneinander im Gesichtsfelde des Fernrohres sieht, zu urtheilen hat, und Freund \textsc{Gerling} die Hauptideen mitgetheilt. Ich hoffe, er wird sein grosses practisches Talent dazu anwenden, diese Idee auszuführen. Übrigens wissen Sie, dass (auf meine Veranlassung) die Aufgabe noch einmal wiederholt ist. . . . . . . . . . . . . . . . . . . . .

Ich komme noch einmal auf das Photometrische zurück. Ich vermuthe, dass das Urtheil über gleiche? oder ungleiche? Helligkeit einer ziemlich grossen Schärfe fähig sein wird, wenigstens einer grösseren, als irgend ein anderes Verfahren, sobald die Gegenstände nicht sehr ungleiche Färbung haben. Im entgegengesetzten Fall, z. B. wenn roth mit grün verglichen werden soll, wird dies freilich wegfallen. Ich habe indess eigentlich keinen recht deutlichen Begriff, \textls[120]{was} eigentlich gleiche Helligkeit hier bedeuten soll. Bei gleichen Farben schaut man das unmittelbar an; aber ungleiche können als Empfindungen in Rücksicht auf Gleichheit der Helligkeit im Grunde kaum verglichen werden. Am Ende bleibt da wohl gar nichts anderes übrig, als dass ein rother Stern und ein grüner Stern für gleich hell gelten, wenn einerlei aliquoter Theil die Grenze des Empfindbaren ist. Etwas ähnliches gilt dann von der specifischen Helligkeit ungleich gefärbter Flächen, wo das Princip auch schon von \textsc{Newton} gebraucht ist. Bei der astronomischen Anwendung möchte indessen daran nicht \textls[120]{so sehr viel} liegen, da doch die sehr abweichend gefärbten Sterne nur seltenere Ausnahmen sind. Aber darin mag wol viel subjectives liegen. Vielleicht wären Augen möglich, wenn auch nicht menschliche, für die das violette Licht (nach jenem Begriff) viel heller wäre als das gelbe. Doch verzeihen Sie, lieber \textsc{Olbers}, meine unreifen Zweifel, da Sie selbst aus diesen Gegenständen viel mehr Studium gemacht\,\raisebox{.15ex}{.}\ haben als ich. . . . . . . .
\clearpage\endgroup
